\documentclass[12pt, a4paper]{article}
\usepackage{amsmath, amsthm, amssymb}
\usepackage{geometry}
\geometry{margin=1in}

\title{Informal Proof of Problem 11.26 (using the Extreme Value Theorem)}
\author{}
\date{}

\begin{document}

\maketitle

\section*{Problem Statement}
Suppose $f$ is a polynomial function of degree $n$ with $f(x) \ge 0$ for all $x \in \mathbb{R}$ (so $n$ must be even). Prove that:
\[
f(x) + f'(x) + f''(x) + \dots + f^{(n)}(x) \ge 0 \quad \text{for all } x \in \mathbb{R}.
\]

\section*{Proof}

Let us define the function $g(x)$ as the sum of $f(x)$ and all of its derivatives up to the $n$-th derivative:
\[
g(x) = \sum_{i=0}^n f^{(i)}(x) = f(x) + f'(x) + f''(x) + \dots + f^{(n)}(x).
\]
Our goal is to show that $g(x) \ge 0$ for all $x \in \mathbb{R}$.

\subsection*{Step 1: The Identity $g(x) - g'(x) = f(x)$}
First, we compute the derivative of $g(x)$:
\[
g'(x) = \sum_{i=1}^{n+1} f^{(i)}(x) = f'(x) + f''(x) + \dots + f^{(n)}(x) + f^{(n+1)}(x).
\]
Since $f(x)$ is a polynomial of degree $n$, its $(n+1)$-th derivative is exactly zero for all $x$, i.e., $f^{(n+1)}(x) = 0$. Therefore, the derivative of $g$ simplifies to:
\[
g'(x) = f'(x) + f''(x) + \dots + f^{(n)}(x).
\]
If we subtract $g'(x)$ from $g(x)$, all the derivative terms perfectly cancel out, leaving only the original function $f(x)$:
\[
g(x) - g'(x) = f(x).
\]

\subsection*{Step 2: Limit Behavior to Infinity}
We are given that $f(x) \ge 0$ for all $x \in \mathbb{R}$. For a polynomial to be non-negative everywhere, two things must be true about its highest-degree term $a_n x^n$:
\begin{enumerate}
    \item The degree $n$ must be \textbf{even}. (If $n$ were odd, the polynomial would approach $-\infty$ on one side of the $x$-axis).
    \item The leading coefficient $a_n$ must be \textbf{strictly positive} ($a_n > 0$).
\end{enumerate}
Now, consider $g(x)$. Taking the derivative of a polynomial decreases its degree. Thus, the polynomials $f'(x), f''(x), \dots, f^{(n)}(x)$ all have strictly lower degrees than $n$.

Consequently, the leading term of $g(x)$ is exactly the leading term of $f(x)$, which is $a_n x^n$. Because $g(x)$ is a polynomial of even degree $n$ with a positive leading coefficient $a_n > 0$, we have:
\[
\lim_{x \to \infty} g(x) = +\infty \quad \text{and} \quad \lim_{x \to -\infty} g(x) = +\infty.
\]

\subsection*{Step 3: The Extreme Value Theorem}
Because $g(x)$ is continuous everywhere and approaches $+\infty$ at both infinities, it cannot decrease infinitely. By a direct consequence of the Extreme Value Theorem on the extended real line, $g(x)$ must attain an \textbf{absolute global minimum} at some point in $\mathbb{R}$. Let us call this point $x_0$.

\subsection*{Step 4: Fermat's Theorem on Stationary Points}
Because $g(x)$ is a polynomial, it is differentiable everywhere. Since $x_0$ is an absolute global minimum of a differentiable function on an open interval (all of $\mathbb{R}$), its derivative at that point must be zero:
\[
g'(x_0) = 0.
\]

\subsection*{Step 5: Completion of the Proof}
We now return to the identity established in Step 1 and evaluate it at our global minimum $x_0$:
\[
g(x_0) - g'(x_0) = f(x_0).
\]
Since $g'(x_0) = 0$, this simplifies to:
\[
g(x_0) = f(x_0).
\]
By our initial hypothesis, $f(x) \ge 0$ for all $x \in \mathbb{R}$. Therefore, we must have $f(x_0) \ge 0$. This immediately implies that:
\[
g(x_0) \ge 0.
\]
Recall that $x_0$ was the \textbf{absolute global minimum} of $g(x)$. This means that for any real number $x$, the value of $g(x)$ is at least as large as $g(x_0)$:
\[
g(x) \ge g(x_0).
\]
Since $g(x_0) \ge 0$, we conclude that $g(x) \ge 0$ for all $x \in \mathbb{R}$. This completes the proof. \qed

\end{document}
